\section{Methodology}
\label{sec:methodology}

We frame database configuration tuning as a black-box optimization problem
over a high-dimensional, mixed-type search space.  The goal is to discover a
PostgreSQL configuration---a vector of runtime parameters called
\emph{knobs}---that maximizes a composite performance score under a given
workload.  Because the relationship between knob settings and database
performance is non-convex, noisy, and riddled with interactions that no
closed-form model captures, we treat the database as an opaque function and
optimize it through repeated evaluation.

This section describes the design of \sysname from four perspectives.  We
first define the configuration space and explain how we select, bound, and
sample knobs (\S\ref{ssec:config-space}).  We then present the
population-based optimization loop and argue why it fits this problem better
than sequential Bayesian alternatives (\S\ref{ssec:pbt-algorithm}).  Next, we
detail the evaluation pipeline---how each candidate configuration is applied,
measured, and scored (\S\ref{ssec:evaluation}).  Finally, we describe the
hardware-aware normalization layer that enables transfer of tuned
configurations across machines with different resource capacities
(\S\ref{ssec:hardware-transfer}).

%% --------------------------------------------------------------------------
\subsection{Configuration Space}
\label{ssec:config-space}

\paragraph{Knob policy and search space definition.}
PostgreSQL~16 exposes roughly 350 runtime parameters.  Most of them---locale
strings, file-system paths, security policies, replication topology---define
\emph{what the database is} rather than \emph{how fast it runs}, and have no
business in an automated tuning loop.  We enforce a formal knob policy that
classifies every parameter into one of eighteen exclusion categories (e.g.,
\texttt{data\_integrity}, \texttt{security\_transport},
\texttt{domain\_scope\_non\_performance}) or marks it as eligible for
optimization.  The policy is recorded in a machine-readable JSON file and
applied at pipeline startup so that every downstream component---sampling,
perturbation, application---operates on the same vetted subset.  Parameters
whose PostgreSQL-native maximum exceeds $2 \times 10^{9}$ (the
\texttt{INT\_MAX} sentinel class) are excluded automatically unless a curated
metadata entry provides practical bounds.  This two-gate design---explicit
exclusion list plus bounds safety gate---prevents the optimizer from
accidentally exploring configurations that corrupt data, crash the server, or
break the tuner's own control channel.

\paragraph{Tiered knob sets.}
Within the eligible parameters, we organize knobs into three nested tiers of
increasing breadth:

\begin{itemize}
  \item \textbf{Minimal (5 knobs):} \texttt{shared\_buffers},
    \texttt{effective\_cache\_size}, \texttt{work\_mem},
    \texttt{random\_page\_cost}, and
    \texttt{max\_parallel\_workers\_per\_gather}.  These are the
    highest-impact parameters; they affect query planning and memory
    allocation for virtually every workload.
  \item \textbf{Core (13 knobs):} Adds checkpoint timing, WAL buffering,
    I/O concurrency, planner statistics targets, and additional parallelism
    controls.  This tier covers the standard production-tuning surface.
  \item \textbf{Standard (${\sim}30$ knobs):} Extends Core with autovacuum
    scheduling, background-writer pacing, WAL-level durability trade-offs,
    and remaining planner cost parameters.  This tier targets comprehensive
    research experiments.
\end{itemize}

The tier system lets practitioners trade search-space dimensionality for
wall-clock budget: a minimal run finishes in minutes on a laptop; a standard
run over thirty knobs may require hours but covers the full tunable surface.

\paragraph{Knob types and sampling.}
Each knob carries a typed definition: \textsc{integer}, \textsc{real},
\textsc{boolean}, or \textsc{enum}, along with curated minimum and maximum
bounds and a scale indicator (\textsc{linear} or \textsc{log}).  We sample
initial configurations from this space using scale-aware random sampling.
For memory and buffer parameters---where the performance impact of doubling a
value from 128\,MB to 256\,MB far exceeds the impact of adding 1\,MB---we
sample in log-space:
\begin{equation}
  v = \exp\!\bigl(\mathcal{U}(\ln v_{\min},\; \ln v_{\max})\bigr),
  \label{eq:log-sample}
\end{equation}
where $\mathcal{U}$ denotes the continuous uniform distribution.  Linear-scale
knobs (e.g., planner cost estimates) are sampled uniformly over their native
range, booleans are sampled with equal probability, and enum knobs are drawn
uniformly from their valid value set.

\paragraph{Two-layer validation.}
We validate configurations at two levels.  The \emph{design-time} layer
checks values against the curated tuning bounds stored in knob metadata---for
example, \texttt{work\_mem} is bounded between 4\,MB and 256\,MB, a range
that reflects practical operating experience.  The \emph{runtime} layer
queries the live PostgreSQL catalog (\texttt{pg\_settings}) for the server's
hard minimum and maximum and rejects any value outside that range.  A
perturbation that overshoots the tuning bound is clamped back; a value that
somehow violates the server's hard bound is caught before the
\texttt{ALTER~SYSTEM} call reaches the database.  This defense-in-depth
pattern keeps the optimizer exploring a productive region while guaranteeing
that no applied configuration can trigger a PostgreSQL rejection.

%% --------------------------------------------------------------------------
\subsection{Population-Based Training}
\label{ssec:pbt-algorithm}

\paragraph{Why PBT over Bayesian optimization?}
The dominant paradigm in recent database auto-tuning work is sequential
Bayesian optimization (BO): fit a surrogate model to observed evaluations,
select the next configuration by maximizing an acquisition function, evaluate,
and repeat~\cite{vanaken2017ottertune}.  BO is highly sample-efficient---it
can reach good configurations in fewer evaluations than random search---but it
is \emph{inherently sequential}.  Each iteration requires fitting the
surrogate, solving the acquisition problem, evaluating one configuration, and
updating the model before the next iteration can begin.

We make a different bet.  Modern infrastructure makes it straightforward to
run $N$ isolated PostgreSQL instances in parallel, each on a separate port.
In this regime, the bottleneck is not the number of evaluations but the
\emph{wall-clock time per generation}: $N$ configurations evaluated
simultaneously cost roughly the same wall time as one.  Population-Based
Training~\cite{jaderberg2017pbt} is designed for exactly this setting.  It
maintains a population of candidate solutions, evaluates them in parallel,
replaces poor performers with perturbed copies of strong performers, and
repeats.  It does not build a surrogate model and therefore avoids the cubic
scaling of Gaussian process fitting.  It also avoids the instability and
sample inefficiency that plague deep reinforcement learning
approaches~\cite{zhang2019cdbtune}, because the worst a worker can do is
score poorly---at which point the exploit step overwrites its broken
configuration with a proven one.

\paragraph{System architecture.}
\sysname organizes the PBT loop around three components.  A
\textbf{Worker}\footnote{Not to be confused with PostgreSQL background
workers.} represents a single population member: it holds the current knob
configuration, the most recent performance score, a step counter, and lineage
metadata recording which elite it last copied from and in which generation.
The \textbf{Evolution} module provides stateless algorithms for exploit and
explore.  The \textbf{Population} orchestrator manages the worker pool,
coordinates parallel evaluation, triggers exploit-explore at the right time,
and detects convergence.

\paragraph{Initialization.}
We create $N$ workers, each holding an independently sampled random
configuration drawn from the knob space described in
\S\ref{ssec:config-space}.  This random initialization seeds the population
with diverse starting points scattered across the search space.  When
warm-starting from a prior session (\S\ref{ssec:hardware-transfer}), we seed
half the population from the archived configuration (with graduated
perturbation at $\pm 20\%$, $\pm 35\%$, and $\pm 50\%$) and fill the
remaining slots with Latin Hypercube Sampling (LHS) to preserve exploration
coverage.

\paragraph{Evaluation.}
At each generation, every worker is evaluated: its configuration is applied to
a dedicated PostgreSQL instance, the target workload is executed, and a
composite score is computed (detailed in \S\ref{ssec:evaluation}).  We run
evaluations in parallel using a thread pool, one thread per worker.  Because
each worker operates on its own database instance (ports 5440, 5441, \ldots),
there is no shared state during evaluation, and thread-level concurrency is
sufficient for the I/O-bound workload execution.

\paragraph{Exploit: truncation selection.}
After all workers in a generation have been scored, we rank them by
performance.  We partition the ranked list into three bands using a quantile
threshold~$\alpha$ (default $\alpha = 0.25$):

\begin{itemize}
  \item \textbf{Elite} (top $\alpha N$ workers): The best-performing
    configurations in this generation.
  \item \textbf{Middle} ($N - 2\alpha N$ workers): Configurations that
    performed adequately and are left unchanged.
  \item \textbf{Poor} (bottom $\alpha N$ workers): Configurations that
    under-performed.
\end{itemize}

Each poor worker copies the full knob vector from a randomly chosen elite
worker.  We pair poor workers with elites uniformly at random rather than
always copying the single best, because deterministic pairing would collapse
population diversity onto one configuration within a few generations.  After
exploitation, the copied worker's step counter resets to zero so that the new
configuration must prove itself through several evaluations before it can be
exploited again.

\paragraph{Explore: perturbation.}
Immediately after exploitation, each copied worker undergoes perturbation.
For every numeric knob~$k$ in the copied configuration, we draw a
multiplicative factor~$\phi_k$ uniformly from $[\phi_{\min}, \phi_{\max}]$
(default $[0.8, 1.2]$) and set
\begin{equation}
  v'_k = \operatorname{clamp}\!\bigl(v_k \cdot \phi_k,\; v_k^{\min},\; v_k^{\max}\bigr),
  \label{eq:perturb}
\end{equation}
where $v_k^{\min}$ and $v_k^{\max}$ are the curated tuning bounds from the
knob space.  Boolean and enum knobs are left unchanged, because toggling them
randomly injects noise without directional signal.  Perturbation ensures that
even when multiple poor workers copy the same elite, their resulting
configurations differ, maintaining the population's diversity around
high-performing regions.

\paragraph{Ready interval.}
A na\"ive implementation that runs exploit-explore after every single
evaluation risks premature convergence: a worker that received a bad score due
to measurement noise could be immediately overwritten.  We adopt the PBT
paper's \emph{ready interval}~$\tau$: a worker must accumulate at
least~$\tau$ evaluations (default $\tau = 3$) before it is eligible for
exploitation.  During those $\tau$ steps, its performance is observed but it
cannot be displaced.  Freshly exploited workers have their step counter reset
to zero, so they too must survive $\tau$ rounds before the next cull.  This
mechanism dampens noisy oscillations and gives newly injected configurations
time to reach steady-state performance.

\paragraph{Stopping conditions.}
Training terminates when any of three conditions is met:
\begin{enumerate}
  \item \textbf{Generation limit:} The population has completed
    $G_{\max}$ generations (a hard safety bound).
  \item \textbf{Early stopping:} The best score across the population has
    not improved for $P$ consecutive generations (default $P = 10$).
  \item \textbf{Convergence:} The standard deviation of worker scores
    falls below a threshold~$\epsilon$ (default $\epsilon = 0.05$),
    indicating that the population has collapsed onto a narrow region and
    further exploration is unlikely to yield improvement.
\end{enumerate}

Algorithm~\ref{alg:pbt} summarizes the complete loop.

\begin{figure}[t]
\begin{minipage}{\columnwidth}
\small
\begin{algorithmic}[1]
  \Require{Knob space $\mathcal{K}$, population size $N$, quantile
    $\alpha$, perturbation range $[\phi_{\min}, \phi_{\max}]$, ready
    interval $\tau$, max generations $G_{\max}$, patience $P$,
    convergence threshold $\epsilon$}
  \State Initialize workers $\{w_1, \ldots, w_N\}$ with random configs
    from $\mathcal{K}$
  \For{$g = 1$ \textbf{to} $G_{\max}$}
    \ForAll{$w_i \in \{w_1, \ldots, w_N\}$ \textbf{in parallel}}
      \State Apply $w_i.\mathit{config}$ to PostgreSQL instance~$i$
      \State Execute workload, collect metrics $m_i$
      \State $w_i.\mathit{score} \gets \textsc{Score}(m_i)$
      \State $w_i.\mathit{steps} \gets w_i.\mathit{steps} + 1$
    \EndFor
    \State Rank workers by score (descending)
    \State $\mathit{elite} \gets \text{top } \lfloor \alpha N \rfloor$
      workers
    \State $\mathit{poor} \gets \text{bottom } \lfloor \alpha N \rfloor$
      workers with $\mathit{steps} \geq \tau$
    \ForAll{$w_p \in \mathit{poor}$}
      \State $w_e \gets$ random selection from $\mathit{elite}$
        \Comment{Exploit}
      \State $w_p.\mathit{config} \gets \textsc{Copy}(w_e.\mathit{config})$
      \State $w_p.\mathit{steps} \gets 0$
      \ForAll{numeric knob $k$ in $w_p.\mathit{config}$}
        \Comment{Explore}
        \State $\phi_k \sim \mathcal{U}(\phi_{\min}, \phi_{\max})$
        \State $w_p.\mathit{config}[k] \gets
          \operatorname{clamp}(w_p.\mathit{config}[k] \cdot \phi_k)$
      \EndFor
    \EndFor
    \If{stopping conditions met}
      \State \textbf{break}
    \EndIf
  \EndFor
  \State \Return best configuration across all generations
\end{algorithmic}
\end{minipage}
\captionof{algorithm}{The \sysname population-based training loop.}
\label{alg:pbt}
\end{figure}

%% --------------------------------------------------------------------------
\subsection{Evaluation Pipeline}
\label{ssec:evaluation}

Each time a worker is evaluated, \sysname must apply the worker's knob
configuration to a live PostgreSQL instance, execute the target workload,
collect performance measurements, and compress those measurements into a
single scalar score.  This subsection describes each stage.

\subsubsection{Configuration Application}
\label{sssec:config-apply}

We apply configurations through a \emph{KnobApplicator} that wraps
PostgreSQL's \texttt{ALTER SYSTEM} interface.  The applicator first queries
\texttt{pg\_settings} to retrieve each parameter's type, valid range, and
\emph{context}---the PostgreSQL term for the scope at which a parameter takes
effect.  Parameters with context \texttt{sighup} take effect after a
\texttt{pg\_reload\_conf()} call; parameters with context
\texttt{postmaster} require a full server restart.

For each knob in the candidate configuration, the applicator validates the
proposed value against the server's hard constraints and then executes a
parameterized \texttt{ALTER SYSTEM SET} statement.  If any parameter fails
validation or application, the entire transaction is rolled back so that the
database is never left in a partially configured state.  After all parameters
are written, the applicator issues \texttt{pg\_reload\_conf()} for
\texttt{sighup}-context parameters and, when necessary, triggers a controlled
restart for \texttt{postmaster}-context parameters such as
\texttt{shared\_buffers}.

A mandatory cooldown period (default 5\,s) follows every configuration change
to let PostgreSQL stabilize---connection pools adjust, background processes
(autovacuum, checkpoint writer) settle, and the buffer pool begins filling
under the new allocation.  Only after cooldown does workload execution begin.

\subsubsection{Workload Execution}
\label{sssec:workload-exec}

\sysname supports three families of workload drivers, all implementing a
common \textsc{SchemaProvider} interface (\texttt{prepare},
\texttt{validate}, \texttt{execute}):

\begin{enumerate}
  \item \textbf{Sysbench (OLTP).}  We delegate workload generation to the
    external \texttt{sysbench} C binary (version~1.1.0+) to eliminate Python
    client overhead.  Supported modes include \texttt{oltp\_read\_only},
    \texttt{oltp\_read\_write}, and \texttt{oltp\_write\_only}, with
    configurable table count, table size, and thread concurrency.  The
    evaluator collects transactions per second (TPS) and p95/p99 transaction
    latency from Sysbench's structured output.
  \item \textbf{TPC-H (OLAP).}  We generate data with the standard
    \texttt{dbgen} binary at a configurable scale factor (default SF\,=\,1,
    producing ${\sim}1$\,GB across eight tables) and bulk-load it via
    PostgreSQL \texttt{COPY}.  We execute the 22~TPC-H decision-support
    queries through \texttt{psycopg2}.  Python client overhead is negligible
    here because individual query execution times range from seconds to
    minutes, dwarfing the millisecond cost of query submission.
  \item \textbf{Custom JSON templates.}  For application-specific tuning,
    users supply a JSON file listing SQL statements with relative frequency
    weights.  The evaluator resolves placeholder tokens
    (\texttt{\{table\}}, \texttt{\{id\}}) against the declared schema and
    executes queries at the specified mix ratio.
\end{enumerate}

Before every evaluation, the executor validates that the benchmark schema
matches the expected configuration (table count, row count, data integrity).
A mismatch triggers automatic schema re-initialization from the snapshot
baseline, ensuring that each worker starts from an identical data state.

A configurable warmup phase (default 100 queries or 10--30\,s) precedes the
measurement window (default 60\,s) to fill the buffer pool and bring the
database to steady state.

\subsubsection{Metric Collection}
\label{sssec:metrics}

We collect three categories of metrics during the measurement window:

\begin{itemize}
  \item \textbf{Latency:} Percentile response times---p50, p95, and
    p99---computed from the full distribution of individual query latencies.
    Percentile metrics are preferable to the mean because a handful of slow
    queries can inflate the average while leaving most users unaffected.
  \item \textbf{Throughput:} Queries per second (QPS) for OLAP workloads;
    transactions per second (TPS) for OLTP workloads.
  \item \textbf{Resource utilization:} CPU utilization, memory utilization,
    disk I/O volume, and buffer cache hit ratio.  We collect these through
    \texttt{psutil}, which reads OS-level accounting directly rather than
    relying on PostgreSQL's internal statistics views.  Accurate resource
    metrics are necessary for the scoring function to distinguish
    configurations that achieve the same throughput but differ in resource
    cost.
\end{itemize}

All metrics are stored in a typed \texttt{PerformanceMetrics} dataclass
that enforces field names, units, and value domains at construction time.

\subsubsection{Scoring Function}
\label{sssec:scoring}

PBT requires a single scalar to rank workers.  We compress the
multi-dimensional metric vector into a composite score through three stages:
normalization, weighting, and gating.

\paragraph{Normalization.}
Raw metrics span incomparable scales (latency in milliseconds, throughput in
thousands of queries per second, utilization as a fraction).  We map each
metric to a utility value $u_i \in [0,\,1]$ using quantile-anchored
normalization.  For metrics where lower is better (latency, error rate), we
invert the mapping.  The normalizer uses the 5th and 95th percentiles of
historically observed values as anchors instead of the raw minimum and
maximum, which makes it resistant to outliers from one-off failures or system
hiccups.  When the fraction of observations falling outside the calibration
bounds exceeds a drift threshold, the normalizer recalibrates from its
retained history and, if many observations saturate the bounds, expands the
anchors to restore discrimination.

\paragraph{Weighting.}
\sysname supports two scoring policies.  The \emph{fixed} policy
(\texttt{fixed\_v1}) assigns static, workload-type-specific weights---for
instance, an OLTP profile that allocates 50\% weight to latency, 30\% to
throughput, 10\% to memory efficiency, and 10\% to error
penalty.  The \emph{feature-driven} policy
(\texttt{feature\_driven\_v2}) computes weights dynamically from a workload
feature vector extracted at runtime.

Under the feature-driven policy, we extract a vector of workload
characteristics---read/write ratio, OLAP complexity, join and aggregation
intensity, concurrency pressure, working-set size, query-mix entropy, and
tail-latency sensitivity---from the benchmark definition.  A linear model
converts these features into per-metric logits:
\begin{equation}
  z_i = b_i + \sum_{j} M_{ij}\, f_j\,,
  \label{eq:logits}
\end{equation}
where $b_i$ is a base bias for metric~$i$, $M_{ij}$ is the feature
coefficient matrix, and $f_j$ is the $j$-th workload feature (with
logarithmic compression applied to working-set size).  We then apply a
floor-constrained softmax:
\begin{equation}
  w_i = \alpha_i + \Bigl(1 - \textstyle\sum_k \alpha_k\Bigr)\;
    \mathrm{softmax}\!\Bigl(\frac{z_i}{T}\Bigr),
  \label{eq:weights}
\end{equation}
where $\alpha_i$ is the minimum floor for metric~$i$ and $T$ is a softmax
temperature.  The floor set satisfies $\sum_i \alpha_i < 1$, guaranteeing
that every metric retains a nonzero contribution while allowing the remaining
weight mass to shift toward the metrics most relevant to the current workload.

\paragraph{Reliability gate.}
Before aggregation, we apply a reliability gate~$G \in [0,\,1]$ that
suppresses scores from unstable evaluations.  If the evaluation suffered a
fatal failure, $G = 0$.  If the error rate exceeded a configurable fatal
threshold, $G = 0$.  For non-zero error rates below the threshold, $G$ decays
linearly toward zero.  Clean evaluations keep $G = 1$.  The gate prevents a
worker that crashed mid-evaluation from receiving a nontrivial score on the
strength of incidental utility values recorded before the crash.

\paragraph{Composite score.}
The final score combines these components:
\begin{equation}
  S = G \cdot \sum_{i=1}^{n} w_i \cdot u_i\,,
  \label{eq:score}
\end{equation}
where the weights $w_i$ sum to one and utilities $u_i \in [0,\,1]$.  The
score is therefore bounded in $[0,\,1]$, gated by reliability, and
interpretable as a weighted average of normalized metric utilities.

%% --------------------------------------------------------------------------
\subsection{Hardware-Aware Normalization and Warm-Starting}
\label{ssec:hardware-transfer}

A configuration tuned on a 2\,GB container is not directly usable on a
32\,GB bare-metal server: absolute byte counts for
\texttt{shared\_buffers} or \texttt{work\_mem} make no sense on a machine
with different physical resources.  To enable configuration transfer across
heterogeneous hardware, we introduce a fractional normalization layer.

\paragraph{Resource detection.}
At startup, \sysname probes the host's available resources---CPU core count
(accounting for \texttt{cgroups} limits in containers), total system memory
(respecting Docker \texttt{--memory} caps), and disk type
(SSD/NVMe~vs.~HDD, detected via rotational flags).  The tuner reserves a 20\%
resource headroom for operating-system processes, setting the effective tuning
budget to 80\% of the detected capacity.

\paragraph{Fractional representation.}
Thirteen knobs are tagged as \emph{hardware-relative}:

\begin{itemize}
  \item \textbf{RAM-dependent:} \texttt{shared\_buffers},
    \texttt{effective\_cache\_size}, \texttt{work\_mem},
    \texttt{maintenance\_work\_mem}, \texttt{temp\_buffers},
    \texttt{wal\_buffers}.
  \item \textbf{CPU-dependent:} \texttt{max\_worker\_processes},
    \texttt{max\_parallel\_workers},
    \texttt{max\_parallel\_maintenance\_workers}.
  \item \textbf{Disk-dependent:} \texttt{random\_page\_cost},
    \texttt{seq\_page\_cost}, \texttt{effective\_io\_concurrency}.
\end{itemize}

When a tuning session concludes, we serialize these knobs as fractions of the
detected resource limits rather than as absolute values.  When warm-starting
on a different machine, the fractions are resolved against the new host's
resources, automatically scaling the configuration to the available hardware.

\paragraph{Aggregate memory validation.}
An optimizer that tunes memory knobs independently can produce configurations
where the combined allocation exceeds physical memory---for instance, by
simultaneously maximizing both \texttt{shared\_buffers} and
\texttt{work\_mem}.  We guard against this with a budget repair pass that
checks the aggregate:
\begin{equation}
  M_{\text{total}} = \texttt{shared\_buffers}
    + \texttt{max\_connections} \times \texttt{work\_mem}
    + \texttt{maintenance\_work\_mem}\,.
  \label{eq:mem-budget}
\end{equation}
If $M_{\text{total}}$ exceeds the per-worker memory budget (80\% of the
node's RAM divided by the number of workers), all memory knobs are scaled
down by a uniform factor $M_{\text{budget}} / M_{\text{total}}$, preserving
the dimensional ratios that the optimizer discovered while enforcing a hard
memory ceiling.

\paragraph{Warm-starting.}
Given a previously saved \texttt{best\_config.json}, \sysname seeds the
population as follows.  One worker receives the archived configuration
verbatim (resolved to the new hardware).  Additional workers receive
graduated perturbations at $\pm 20\%$, $\pm 35\%$, and $\pm 50\%$, up to
$N/2$ workers total.  The remaining $N/2$ slots are filled by Latin Hypercube
Sampling over the full knob space.  This split between exploitation of prior
knowledge and fresh exploration prevents negative transfer: if the archived
configuration is poorly suited to the current workload, the LHS-seeded
workers provide a fallback, and the PBT exploit step will displace the
warm-started workers within a few generations.
